\section{Pilot Study}
\label{sec:pilot}

We ran three pilots before the full study: one in each condition (sgt and git,
both on the same task set) and a third exercising the study apparatus end to end.
The first two were designed to find defects in the tool that would cost a real
participant their data point, and defects in the study design that would make a
result uninterpretable. The third tested the machinery between the participant's
machine and the facilitator's console---exactly the seam nobody had
exercised---and found a validity threat that neither task pilot could surface.
We report all three because the defects they found are consequences of design
commitments rather than accidents of implementation: a technical evaluation
establishes that a property holds across 10{,}237 operations; it takes a new
user to show that a property that holds still costs too much to use.

\subsection{What the pilots tested}

Six requests, drawn from the scenarios of Section~\ref{sec:scenarios}: find
which commit introduced two unrelated changes (provenance), remove a named
feature while keeping the rest (removal), selectively restore a single function
(selective restore), locate and repair a regression whose test still passes
(regression repair), plan work on a branch and resolve a conflict (plan and
fork), and split a tangled commit so each half has a clear name (history edit).
Both conditions worked on the same five-thousand-line Python project and were
scored with the same rubric and the same automated check.
Table~\ref{tab:pilot} summarises the outcomes.

\begin{table}[t]
\centering
\small
\begin{tabular}{p{2.8cm}p{2.4cm}p{2.2cm}}
\toprule
Request & \sgt{} condition & git condition \\
\midrule
R1: Provenance & Correct, 4 cmds & Correct, 6 cmds \\
R2: Removal & Done; tool corrupted file, repaired by hand & Done cleanly \\
R3: Selective restore & \cmd{restore} crashed; recovered via git & Done cleanly \\
R4: Regression repair & Correct fix + better test & Correct fix + better test \\
R5: Plan \& fork & Done in git (declined sgt) & Done cleanly \\
R6: History edit & Half done (no rename verb) & Done (scripted rebase) \\
\bottomrule
\end{tabular}
\caption{Per-request outcomes across both pilot conditions on coursecraft (same
tasks, same rubric). The sgt condition ran on version 0.1.0; twelve defects
fixed afterwards. R4 and R5 produced equivalent results in both conditions.}
\label{tab:pilot}
\end{table}

\subsection{What the sgt condition found}
\label{sec:pilot-sgt}

The sgt pilot completed four of six tasks and reached the correct answer on
a fifth. Three of the eight verbs in the daily interface failed during the
session: \cmd{revert} produced syntactically invalid Python, \cmd{restore}
crashed with a raw traceback, and \cmd{save} refused to record a legitimate
edit. The participant fell back to git for the repairs and finished the
substantive work outside the tool.

Twelve defects were fixed from this one session. We describe three because they
are consequences of design commitments this paper defends.

\paragraph{Semantic subtraction glued lines together.}
Reverting the waitlist feature left one file with three function headers on one
line, \cmd{find\_student} relocated to end of file, and thirteen collection
errors in pytest---under a green checkmark. The cause is the round-trip
design of Section~\ref{sec:sets}: the rebuild synthesises zero separator bytes,
so a removal that takes the layout chain of an entity it keeps produces
definitions with no whitespace between them. The fix re-grounds the layout of
every kept entity after a subtraction.

\paragraph{A preview flag mutated state.}
The tutorial teaches that every destructive verb shows a preview first and acts
only with \cmd{--yes}. Adding \cmd{--keep-dependents} to a revert silently
voided that contract, writing intermediate state into the store with no preview
and no confirmation. The participant's verdict: ``If a flag can turn a preview
into an action, I can't reason about which of these commands are safe.'' The
flag now requires \cmd{--yes}.

\paragraph{The map showed husk features.}
\cmd{sgt log --map} displayed a feature called \emph{Section Waitlist} that
contained zero symbols in zero files---a cluster whose only records were
bookkeeping sentinels---while the feature that actually held the waitlist
code was named after a deleted experiment and did not appear in the view at all.
The participant followed the wrong name for the entire removal task. Later
validation found the problem was systemic: 9 of 24 leaf features in each study
repository owned no behavioural code, including the three features the study's
own removal requests name. The mechanism was that the clustering graph includes
positional gap-bytes as nodes (they carry within-file cohesion), and Leiden
readily forms a community made entirely of them. The fix absorbs any leaf whose
members are all non-behavioural into the leaf that already owns the entities
those members hang off, reducing coursecraft from 24 leaves to 14 with none
empty.

\smallskip\noindent
The comprehension operations (naming what a request produced, seeing what a
revert would remove) held up under a user who had never seen the tool before;
the participant credited ``sgt's attribution and its revert preview, which git
genuinely cannot do.'' The mutation operations did not survive contact, and the
data-loss path that one of them opened---a \cmd{fulfill} command that overwrote
uncommitted work in five files while printing a green checkmark---was the worst
defect we found in the system during this evaluation. The participant's
conclusion: ``From here I treat \cmd{sgt revert} as `a fast, accurate first
pass that leaves me a punch list', not as an operation that completes.'' When
asked at exit whether they would adopt the tool: ``For reading history, yes,
starting today. For anything that writes, no.''

\subsection{What the git condition found}
\label{sec:pilot-git}

The git pilot completed all six tasks inside the time budgets, including
request~6 (history edit), which the protocol had expected to be out of reach in
the git condition. The participant scripted the interactive rebase, split the
tangled commit in two, and verified that both halves pass independently by
comparing tree hashes. That result is better than the sgt condition managed on
the same task: sgt had already split the same commit at mining time, which is
the more impressive half, but the participant could not give the two halves
clear names because no verb renames a checkpoint.

Four defects in the study design surfaced:
\begin{enumerate}
\item The git-condition copies carried commits written by \sgt{}
  (\cmd{sgt revert f-10462e17@2}, \cmd{sgt undo: restore prior ideal}).
  A participant in the baseline should not be able to tell that another tool was
  involved.
\item A request ticket described a change as ``showing times differently'' when
  no earlier format existed. Both pilots spent time on our wording rather than
  on the task.
\item Request~6 needs a rubric for both conditions rather than an assumption
  that the git condition cannot finish.
\item One ticket referenced a duplicate commit and an undo authored by the
  researcher.
\end{enumerate}
All four are fixed.

The git participant also described, unprompted, what they wished the history
could answer: ``Ask history a question about a symbol, not a file. `What
happened to \sym{overlaps}, who calls it, and when did that change?' would have
turned request~4 from a hunch into a lookup.'' A baseline participant
independently describing the treatment condition's premise is evidence that the
problem this paper addresses exists independently of the tool we built. The
need exists before the solution is shown, and we have added a question to the
exit interview to collect this evidence deliberately.

\subsection{What the apparatus pilot found}
\label{sec:pilot-apparatus}

The third pilot ran the full study end to end---real bundles, real terminal, real
git and \sgt{}---with the sole purpose of exercising the machinery between the
participant's machine and the facilitator's console. It found nine defects, all
in the seam between systems. Three generalise.

\paragraph{The primary read verb rejected its own output.}
\cmd{sgt show} failed 6 of the 10 times it was used, and it is the verb the
practice sheet tells participants to lean on. It rejected save identifiers printed
by \cmd{sgt log}, full git commit hashes, and feature names typed exactly as
displayed. The cause was narrower than ``identifiers are inconsistent'': the id
column of \cmd{sgt log} prints the 7-character git commit hash, and \cmd{show}
did not accept that format as a selection kind. The participant fell back to
\cmd{git blame} and \cmd{git show --stat}---in the \sgt{} condition---and
answered the question that way. A tool that pushes someone back to the baseline
mid-task is not a measurement problem; it is the measurement. The defect is
fixed: \cmd{show} now resolves commit-shaped tokens as saves.

\paragraph{A validity threat: the sgt condition degrades its own baseline.}
\label{sec:pilot-trailers}
Every commit in the sgt-condition repositories carries its operation identifiers
as commit-message trailers. On a representative commit whose actual message is
the single line \texttt{add course search}, we measured 174 \texttt{Sgt-Op:}
trailer lines. So \cmd{git show}, \cmd{git log}, and \cmd{git log -p} all bury
the real content under a screenful of hex. The git-condition repositories are
stripped of these by design, so a participant in the baseline sees clean history
while a participant in the treatment sees this. The direction of the bias is the
problem: it makes git look worse exactly where \sgt{} is being measured against
it, and it does so in the treatment arm---the arm whose participant reaches for
\cmd{git blame} as a fallback whenever \sgt{} does not have the answer. We
disclose this because it favours our condition: a win for \sgt{} on any task
where the participant consulted git history is weaker than it appears, and we
cannot retrospectively subtract the effect. The pre-registered analysis
acknowledges it as a confound.

\paragraph{Carry-over between isomorphic tasks.}
The two test repositories are isomorphic by design (same episode script, nouns
swapped) so nobody sees the same project twice. The pilot participant reported
that the resemblance is strong: ``by half two I was pattern-matching against an
already-solved puzzle rather than reading fresh.'' The design counterbalances
order and models it as a factor; this finding confirms that modelling order is
essential rather than a formality, and that the carry-over effect is likely large
relative to the condition effect on a twelve-person sample.

\subsection{Thematic analysis}
\label{sec:pilot-lessons}

We coded all three pilots' observations against the design commitments of
Section~\ref{sec:design}, interpreting each through three lenses from the
automation trust literature: Parasuraman et al.'s four-stage model of
automation~\cite{parasuraman2000} (which stage failed?), Lee and See's trust
calibration framework~\cite{lee2004} (what did the failure teach the user about
the system's reliability?), and Norman's gulfs of evaluation and
execution~\cite{norman1986} (could the user perceive the state, and could they
act on it?). Six themes emerged that generalise past the individual defects.

\paragraph{Comprehension held; mutation did not.}
The comprehension operations---attribution, revert preview, dependency
display, feature listing---produced correct answers on first contact in the task
pilot. The sgt participant reached the right answer on request~1
in four commands, using a path the git participant could not take (symbol-level
attribution across files). Every write operation---\cmd{revert},
\cmd{restore}, \cmd{save}, \cmd{fulfill}---either failed, crashed, corrupted
output, or refused to run. The apparatus pilot qualifies this split: the
\emph{interface} to the comprehension layer failed at the same rate as the
mutation layer (\cmd{show} rejected 6 of 10 queries), so a user who could not
navigate to the right screen never reached the comprehension that screen
provides. The distinction matters because it separates two failure modes with
different fixes: the underlying analysis was correct (the attribution existed and
could be displayed), but the path to it was blocked by an identifier mismatch
between two surfaces of the same tool.

Norman calls these two directions the gulf of
evaluation and the gulf of execution: the first asks whether a user can tell
what state the system is in, the second whether they can act to change
it~\cite{norman1986}. Parasuraman, Sheridan, and Wickens decompose automation
into four stages---information acquisition, information analysis, decision
selection, and action implementation---and argue that automating the earlier
stages while leaving the later ones to the human is both safer and more robust
than automating uniformly across all four~\cite{parasuraman2000}. \sgt{}
automates stages one and two heavily (parsing, clustering, dependency tracking,
consequence display) and leaves stage three to the developer (which feature to
revert, how to resolve a fork). The pilot failure sits squarely at stage four:
the action implementation that follows a correct analysis either crashed or
corrupted output. The design implication is that the comprehension layer is
independently valuable even when the action layer is not yet trustworthy, but
only if the developer can reach it---an interface failure at the front door costs
the same as a corruption behind it, because in both cases the developer falls
back to git.

\paragraph{Silent wrong answers erode trust faster than crashes.}
The participant encountered both an omission error (\cmd{restore}: raw traceback,
no data loss, the system failed to act) and a commission error (\cmd{revert}:
green checkmark, three function headers glued onto one line, the system acted
wrongly while reporting success). Mishler and Chen showed that commission errors
erode trust in automation faster than omission errors, because a system that does
the wrong thing while claiming to have done the right thing attacks the
operator's ability to distinguish reliable reports from unreliable
ones~\cite{mishler2023}. That asymmetry appeared here in full: the crash cost one command; the
corruption cost twenty minutes of repair work, left the participant unwilling to
trust checkmarks, and propagated forward: when the session feature later
reported ``86 new ops'' on a workspace with zero edits, the participant
declined to use it. The number was explainable, but the tool had spent its
credibility on an earlier false~positive. Lee and See's model predicts exactly this: a crash provides
correct calibration information (the system failed and said so), whereas a false
positive undermines the calibration mechanism itself, because the user can no
longer distinguish a true positive from the next false
one~\cite{lee2004}. Wang et al.\ found the same asymmetry in developers using
AI code generation: one bad recommendation erases the credibility of three good
ones, and a single negative experience with a category of output causes
developers to stop reading that category
entirely~\cite{wang2024trust}. Madhavan and Wiegmann's integrative review
explains why the recovery never came: people hold automation to a higher initial
standard than they hold humans, and a single failure triggers categorical
attribution---the system is \emph{broken}---with no intermediate state between
working perfectly and being untrustworthy~\cite{madhavan2007}. Sabouri et al.\
measured the downstream consequence: developers retained only 52\% of initially
accepted AI suggestions on revisitation, and the retention rate dropped further
after encountering a single incorrect suggestion they had initially
accepted~\cite{sabouri2025}. Trust is not per-operation but cumulative, and one
self-reporting failure costs more than the task it belongs to.
Eight such failures appeared across the full evaluation period,
four of which we had not found after eight months of daily use and
found only by auditing the evaluation's own results---the newest was a revert
that reported ``removes 0 edits'' and then rewrote two files, a command whose
entire pitch is legible consequences telling the developer nothing will happen
before changing their code.

\paragraph{The tool's value was in the punch list.}
The participant reframed the revert as ``a fast, accurate first pass that
leaves me a punch list'' rather than an operation that completes on its own. Every item on the
removal's consequence report came true: the three functions that still
referenced removed code broke, and nothing else did. The preview was the
product; the execution was the problem. The participant's exit debrief sharpened
this: ``It indexes your history by function instead of by file, so you can ask
`what changed this symbol, and what else came along for the ride' and get a real
answer in one command---that part is excellent. Use it to find things, then make
the change yourself with git.'' Ferdowsi et al.\ found that continuously displaying runtime values mitigates
both over-reliance and under-reliance on AI-generated code, by lowering the cost
of validation through execution~\cite{ferdowsi2024}. The preview serves the same
function: it lowers the cost of evaluating a revert by showing the consequence
before it happens, and the pilot confirms both halves of Ferdowsi's finding.
Where the preview was correct (the three named functions broke, nothing else
did), the participant relied on the tool appropriately. Where it was incorrect
(the green checkmark over corruption), the lowered validation cost produced
worse outcomes than no tool at all, because the participant trusted the preview
and did not verify independently. A partially-working tool can
deliver value through its read layer alone---it tells the developer what to
change, and they change it with the tools they already trust---provided the
reports are accurate. The accuracy condition circles back to the previous theme:
a tool whose punch lists are wrong on one task stops being consulted on the next.

\paragraph{Vocabulary mismatch across surfaces.}
Four counting units---edits, ops, symbols, saves---appeared in different
outputs with no definitions, and the participant stopped reading the numbers
halfway through: ``the edit counts have been meaningless all session.'' Two
views of the same feature reported different symbol counts (the tree counted
cluster members; \cmd{show} counted footprint symbols). Selector namespaces
varied per verb, with no cross-lookup. Each of these is a consistency defect,
not a missing feature, and each one makes the tool harder to learn because the
participant cannot build a stable model of what the numbers mean. The fix was
mechanical: one counting unit (symbols) and one namespace across all verbs.
Carroll and Rosson named the mechanism that makes this costly: active users
learn by generalising from one interaction to the next, and inconsistent surfaces
punish generalisation by teaching a rule that fails on the second
screen~\cite{carroll1987}. Jackson frames the same requirement as a property of
the concepts a system exposes: a concept whose state looks different depending on
which operation queries it has no stable operational
principle~\cite{jackson2021}. The problem compounds in a tool whose labels are
already inferred: Nguyen and Glassman showed that vocabulary mismatch between a
developer and a model is the norm rather than the exception~\cite{nguyen2024},
so a developer already arrives at the tool uncertain of the code's naming; a tool
that then disagrees with itself across its own screens removes the one stable
reference point (the tool's own output) that could have resolved that
uncertainty. The lesson is that internal consistency is a learnability
precondition (not merely cosmetic), and a tool whose surfaces disagree teaches
the user to ignore its numbers.

\paragraph{Baseline sufficiency reframes the study.}
The git condition completed every task, including the one we had expected to be
out of reach (splitting a tangled commit via scripted interactive rebase). The
study therefore measures cost and confidence rather than raw capability: both
conditions can finish, and the question is how many commands, how many errors,
and how much the developer trusts the result. That is a stronger design than one
where the baseline cannot finish, because it makes both directions of the
comparison informative---the git condition can win on tasks where the tool's
mutation failures cost more time than the workaround saves.

\paragraph{The treatment degrades its own fallback.}
The apparatus pilot found that the sgt condition makes plain git---the tool a
participant reaches for when \sgt{} fails---measurably worse than in the git
condition. Every sgt-authored commit carries operation identifiers as
commit-message trailers (174 lines on a one-line commit), so \cmd{git show} and
\cmd{git log} bury the content a developer needs under hex they cannot use. The
git-condition repositories are stripped of these by design, so the baseline sees
clean history while the treatment does not. This creates a confound that favours
our condition: any comprehension task where the participant consulted git history
in the sgt arm was harder than the identical task in the git arm, for a reason
unrelated to \sgt{}'s attribution. The measurement lesson is that a tool layered
on top of a host system (here, git) must account for how it changes that host's
usability, because a participant evaluating the treatment is also evaluating a
degraded version of the control. We acknowledge this as a threat and note that
the pre-registered analysis cannot correct for it post hoc.

\medskip\noindent
The six themes, read together, describe a tool whose analytical core is sound
and whose interface to the developer is not yet stable enough to carry the
analysis through to action. The first three themes (comprehension/mutation
split, trust erosion, punch-list reframing) all resolve in the same direction:
the developer can use this tool as a comprehension instrument today, and should
not trust it as an execution instrument yet. The fourth (vocabulary mismatch)
explains why even the comprehension layer requires polish before study: a
developer who cannot navigate to the right screen never reaches the analysis
that screen provides, and inconsistency between screens teaches them to stop
reading. The fifth (baseline sufficiency) constrains the study design: both
conditions finish, so the measurement must be about confidence and error rate
rather than completion. The sixth (treatment degrades control) constrains the
interpretation: any advantage we observe in the treatment arm is weakened by a
confound we cannot remove.

The twenty-one fixes and four design corrections together mean that the study we
will run differs substantially from the study we would have run two weeks
earlier. Three consequences follow for its design. First, the primary dependent
variable is now time and error count per task rather than task completion, since
both conditions finish. Second, the exit interview asks whether the participant
trusts the tool's reports, because the pilot showed that credibility, once lost
to a false positive, does not recover within a session. Third, the sgt
condition's tasks are scored on the comprehension half (did the participant
identify the right target) separately from the mutation half (did the tool
execute correctly), so that a mutation failure in the tool does not mask a
comprehension success in the participant. The result that would count against the
design is error-rate parity between the two conditions. If both conditions
produce equal rates of undetected errors, then the recovery this paper describes
is something developers already do reliably without the tool, and
Section~\ref{sec:scenarios} overstates the difficulty. We
report the pilots rather than only the findings, because we want a reader to
know which version of the tool the full study runs against and why the study
measures what it measures.

\subsection{Study repository validation}
\label{sec:pilot-repos}

We built two 5{,}000-line Python projects (coursecraft and confplan) using
scripted agent sessions---28 and 26 commits, 23 episodes each---so that the
ground truth for every task is known from the build log. Running the
derivability census on both repositories before the pilot produced one pass and
one failure: confplan's central removal task (``take the promotion system out'')
was not answerable from the record, because the promotion subsystem sat inside
a 40-symbol feature that also held the entire CLI. The same episode script, the
same 23 episodes, the same task---answerable on one repository and not the other.

Tracing the cause exposed a single defect with four visible consequences: the
largest feature was named \emph{init repo} (a label derived from one seed-commit
vote in a 71-op leaf), 77--84\% of commits had edits filed under that seed
feature, eight features carried file-path labels, and confplan's promotion
subsystem was absorbed into one cluster. The mechanism was provenance blindness
on the save path: \sgt{}'s two clustering signals read each operation's
provenance field to learn ``these symbols were written in the same sitting,''
and every operation saved through \cmd{sgt save} carried no provenance---the
commit reference lives in trailers that three call sites did not resolve. 366
of 370 operations, unattributed. Co-change without temporal separation is
co-occurrence, which produces one cluster.

The finding that generalises beyond our system: the blindness appeared only in
histories built \emph{through} \sgt{}. Every unit test, every dogfood run on
the development repository, and every retroactively-mined fixture had
provenance, because all of those derive their temporal signal from git commits
that already exist. The save path---the one workflow we ask participants to
use---was the least-tested path in the codebase, and a tool evaluated only
against mined history has been evaluated on its strongest arm.

After three fixes---provenance resolution, label prompt filtering, and an
aliasing-rename crash---both repositories pass an 8-of-8 derivability gate: the
answer to each of the study's four requests is reachable through the taught verb
set (\cmd{log}, \cmd{show}) in at most three commands. We score this as ``8/8
by the author's reading'': the pre-registered primary metric could not be
computed as written, so a referee is entitled to treat the replacement as weaker
evidence.

Four tangle cases---commits the build log placed two unrelated changes into---are
all separated correctly: in coursecraft, \sym{cli.py::cmd\_search} lands in
one feature while \sym{slots.py::parse\_slot} and its test land in another, and
the two concerns the commit message conflates are two features in the record.
Reconstruction is exact on both repositories (0 drifted paths), so nothing
above came at the cost of fidelity.

One asymmetry remains: coursecraft's waitlist feature contains five enrollment
symbols (a hub function's edit chain pulls them in), so a feature-level revert
would violate the request to ``keep enrolling working.'' The member list is
visible before applying, which makes this a legibility cost rather than a silent
failure, and it is the cost we report to participants in the task handout
rather than one we hide. On mirrored histories the clustering diverges---the same
episode script produces 2 separable features in coursecraft and 3 in
confplan---which is itself a finding about the method's stability under
near-identical inputs.
